% Schematic examples derived from the stated equations; not experiment data.
\begin{figure}[H]
\centering
\begin{minipage}[t]{0.32\textwidth}
\centering
\textbf{(a) Channel probability mass}\par\smallskip
\begin{tikzpicture}[x=1cm,y=1cm,font=\scriptsize]
\path[use as bounding box] (-0.1,-3.8) rectangle (5.1,0.4);
\node[anchor=west] at (0,0) {One softmax: illustrative topic};
\filldraw[fill=accent!70,draw=black!65] (0,-0.65) rectangle (3.74,-0.2);
\filldraw[fill=lossgold!70,draw=black!65] (3.74,-0.65) rectangle (4.4,-0.2);
\node at (1.87,-0.425) {fragments: 0.85};
\node at (4.07,-0.88) {losses: 0.15};
\draw[dataflow] (2.2,-1.05)--(2.2,-1.45);
\node[anchor=west] at (0,-1.8) {Separate normalizers};
\filldraw[fill=accent!70,draw=black!65] (0,-2.45) rectangle (2.2,-2.0);
\filldraw[fill=lossgold!70,draw=black!65] (2.2,-2.45) rectangle (4.4,-2.0);
\node at (1.1,-2.225) {fragments: 0.50};
\node at (3.3,-2.225) {losses: 0.50};
\node[align=left,text width=4.7cm,anchor=north west] at (0,-2.85)
{Topic probability mass is balanced.\\Observed counts are unchanged.};
\end{tikzpicture}
\end{minipage}\hfill
\begin{minipage}[t]{0.32\textwidth}
\centering
\textbf{(b) Whole-spectrum context}\par\smallskip
\begin{tikzpicture}[x=1cm,y=1cm,font=\scriptsize]
\path[use as bounding box] (-0.1,-3.8) rectangle (5.1,0.4);
\node[anchor=west] at (0,0) {Same word: $\widehat\rho_w=(1,0)$};
\draw[black!25] (0.35,-2.65)--(0.35,-0.45);
\draw[black!25] (0.1,-1.55)--(4.4,-1.55);
\draw[dataflow,draw=accent,line width=1pt]
  (0.35,-1.55)--(2.82,-0.727);
\draw[dataflow,draw=lossgold,line width=1pt]
  (0.35,-1.55)--(2.82,-2.373);
\draw[dataflow,line width=1pt] (0.35,-1.55)--(2.95,-1.55);
\node[anchor=west,text=accent] at (2.90,-0.70) {$h_{Aw}$};
\node[anchor=west] at (3.02,-1.55) {$\widehat\rho_w$};
\node[anchor=west,text=lossgold] at (2.90,-2.40) {$h_{Bw}$};
\node[align=left,text width=4.7cm,anchor=north west] at (0,-2.85)
{A: $s_A=(0.5,0.5)$\\B: $s_B=(0.5,-0.5)$\\
$h=\operatorname{normalize}(\widehat\rho_w+s)$};
\end{tikzpicture}
\end{minipage}\hfill
\begin{minipage}[t]{0.32\textwidth}
\centering
\textbf{(c) Sparse topic proportions}\par\smallskip
\begin{tikzpicture}[x=1cm,y=1cm,font=\scriptsize]
\path[use as bounding box] (-0.1,-3.8) rectangle (5.1,0.4);
\node[anchor=west] at (0,0) {Identical scores: $z=(1,0,-1)$};
\node[anchor=west] at (0,-0.4) {Softmax};
\filldraw[fill=accent!70,draw=black!65] (0,-1.05) rectangle (2.927,-0.65);
\filldraw[fill=lossgold!70,draw=black!65] (2.927,-1.05) rectangle (4.004,-0.65);
\filldraw[fill=black!20,draw=black!65] (4.004,-1.05) rectangle (4.4,-0.65);
\node[anchor=west] at (0,-1.3) {$(0.665,\;0.245,\;0.090)$};
\node[anchor=west] at (0,-1.8) {1.5-entmax};
\filldraw[fill=accent!70,draw=black!65] (0,-2.45) rectangle (3.655,-2.05);
\filldraw[fill=lossgold!70,draw=black!65] (3.655,-2.45) rectangle (4.4,-2.05);
\node[align=left,text width=4.7cm,anchor=north west] at (0,-2.85)
{$(0.831,\;0.169,\;0)$\\Topic 3 has exactly zero weight.};
\end{tikzpicture}
\end{minipage}
\caption{Illustrations of the three enhancements, not measured model outcomes.
In (a), bar length is total probability in each spectral channel. In (b),
a shared word vector is tilted by two different spectrum means, with context
strength one. The resulting unit directions are approximately
$h_{Aw}=(0.949,0.316)$ and $h_{Bw}=(0.949,-0.316)$; the axes are embedding
coordinates. In (c), colored segments are topic 1, topic 2 and
topic 3, in that order; each complete bar has unit mass.}
\label{fig:enhancements}
\end{figure}
